\section{How to use the audit}
\label{sec:use}
Suppose a team wants to compare monitors for an order-management agent. The team should first write the question it wants answered: ``Did the agent make the requested database change?'', ``Did it follow the one-call rule?'', or ``Was each refund permitted?'' A single success label cannot answer all three. The practical output of this audit is a comparison of warnings against named checks, with the evidence available at each warning.

\paragraph{Start with one recorded conversation.}
For the illustrative order example in the introduction, record three separate facts: the final order is correct; the two-call proposal breaks the interaction rule; and the monitor warns when it sees that proposal. A warning for the one-call rule is justified. The same record does not demonstrate advance prediction, nor does the correct order establish compliance with every other policy. This small inspection should precede any aggregate error rate.

\paragraph{Then apply the same checks to the full evaluation set.}
\begin{enumerate}
\item \textbf{State the monitor's job and what it can see.} Name the rule or outcome, and specify whether the monitor sees a proposed action before execution or only its result afterward.
\item \textbf{Label that property separately from final success.} Use executable checks when the policy and observed facts make them possible. Keep unknown facts distinct from passed checks. Record the boundary of each check; a passed selected predicate cannot certify requirements outside its coverage.
\item \textbf{Compare monitors on the same evidence.} Save the policy, tool definitions and conversation through each decision. Exclude future turns and outcome labels. Record model settings and thresholds, and include a simple rule checker where one is available.
\item \textbf{Report counts with readable denominators.} For example, write ``two warnings among 49 successful conversations that followed the interaction rule.'' Also report warnings on violating conversations, missing outputs, distinct tasks, repeat runs and when warnings occurred. A group with no examples has an undefined rate.
\item \textbf{Investigate disagreements before changing the monitor.} A warning may expose a bad reference label, a different monitoring target, an unmeasured rule, or a model mistake. If warnings trigger blocking, place the gate before write execution and return its reason to the continuing agent. Include permitted controls, record whether blocked actions would already fail in the native tools, and measure both prevented changes and subsequent conversation outcomes. Keep one-step replay distinct from a counterfactual full conversation.
\end{enumerate}

\paragraph{Reproducing the reported results.}
Saved-record analyses require no paid model calls. The earlier reproduction workflow in \texttt{USAGE.md} and \texttt{reproduce.sh} covers the reserved-test cohorts, cancellation fixtures, pilot/new-task analyses and initial live-gating comparison. It does not reproduce the later compact study or assemble the combined paper. The combined-paper reporting checks and source builder are listed in its artifact-map appendix. The studies remain separately reported: \texttt{replication/analysis.json} holds the reserved-test and cancellation comparisons; \texttt{substantive\_evaluation/analysis.json} holds the later matched-input comparisons; and \texttt{intervention\_study/live\_gating/analysis.json} holds the initial live comparison. The compact study has its own frozen records and analysis. Preserved replay checks and optional fresh native replay are distinguished from fresh model collection.

\paragraph{What this enables.}
The audit helps a team tell whether a reported warning or block rate answers its intended question. The live studies add specific prevented writes and subsequent conversation outcomes; their scope remains the tested configurations and scenarios. This reporting practice is consistent with the NIST AI RMF's emphasis on appropriate metrics and documenting unmeasured risks \citep{nist2023rmf}, and with the generative-AI profile's caution about extrapolating from narrow assessments \citep{nist2024genai}. Application-specific policy coverage and tests of actual interventions remain necessary for operational claims.
